\begin{DeletedBlock}
\PurePy is a subset of Python: every well-formed program is a syntactically valid Python program, and every terminating run with defined behaviour under \PurePy{} produces the same result under Python. \figref{taxonomy} classifies programs and runs by their outcomes under the two
languages. The hatched cells of the figure are empty because every well-formed program is a syntactically valid
Python program, and each result in the last column is the same as under Python.
\end{DeletedBlock}
\begin{AddedBlock}
A Python program is \emph{syntactically valid} if it does not raise \pyinline{SyntaxError} under Python. A Python
program is a \emph{\PurePy program} if it has the syntax of \secref{syntax}, and a \PurePy program is
\emph{well-formed} if it satisfies the well-formedness rules of \secref{well-formedness}. Some of the operations excluded by \PurePy, such as comparison of two functions, cannot be
rejected statically\paperonly{ (\secref{no-coercions})}\speconly{ (summarised in the informative
\annexref{exclusions})}. A run which reaches one of these \emph{dynamically excluded} operations has undefined
behaviour, and aborts with $\errUndef$. \PurePy is a subset of Python: every well-formed program is syntactically
valid, and every terminating run with defined behaviour under \PurePy produces the same result under Python.
\figref{taxonomy} classifies programs and runs by their outcomes under the two languages. The hatched cells of the
figure are empty because every well-formed program is syntactically valid. Each result in the last column is the same
as under Python.
\end{AddedBlock}
